% sections/14_threat_model.tex
\section{Threat Model and Trust Assumptions}\label{sec:threat-model}

This section defines the adversary model, security goals, trust anchors,
and key-management policy governing the \ISI integrity system.  The model
is scoped to the computational pipeline, snapshot artefacts, and serving
layer.  Raw-data veracity is explicitly out of scope.

Normative keywords (SHALL, SHALL NOT, MUST, MUST NOT, MAY, RECOMMENDED)
are used in the sense defined by \textcite{rfc2119} throughout this section.

% -------------------------------------------------------------------------
\subsection{Adversary Classes}\label{sec:threat:adversaries}

The following adversary classes are considered.  Each class is designated
\textbf{A-\textit{n}} and is referenced in the mitigation analysis below.

\begin{description}[style=nextline, leftmargin=3.5em, labelwidth=3em]

  \item[\textbf{A-1}] \textbf{Post-materialisation file tampering.}\enspace
    An adversary modifies one or more files in a published snapshot directory
    after the snapshot has been signed.  The attack surface includes score
    files, the methodology registry, and metadata files.  The adversary does
    not possess the signing key.

  \item[\textbf{A-2}] \textbf{Partial snapshot substitution.}\enspace
    An adversary replaces a subset of files in a snapshot directory
    with files from a different vintage or methodology version while
    leaving the remaining files intact.  The goal is to produce a snapshot
    that passes superficial inspection but contains inconsistent data.

  \item[\textbf{A-3}] \textbf{API response manipulation.}\enspace
    A network-level adversary intercepts and modifies API responses between
    the serving infrastructure and the client.  This includes modification
    of JSON payloads, HTTP headers, and ETag values.

  \item[\textbf{A-4}] \textbf{Registry mutation.}\enspace
    An adversary with write access to the deployment environment modifies
    the \texttt{methodology\_registry.json} to alter parameter values
    (e.g., classification thresholds, rounding precision) without
    recomputing the affected scores.

  \item[\textbf{A-5}] \textbf{Floating-point drift.}\enspace
    A non-conforming implementation produces scores that differ from the
    canonical output due to platform-specific floating-point behaviour
    (extended precision, FMA contraction, non-IEEE rounding mode).  The
    divergence is unintentional but produces hash mismatches and
    potentially different classifications.

  \item[\textbf{A-6}] \textbf{Key compromise.}\enspace
    The Ed25519 signing key is obtained by an adversary, enabling
    production of valid signatures over fabricated snapshots.

\end{description}

% -------------------------------------------------------------------------
\subsection{Security Goals}\label{sec:threat:goals}

The integrity system SHALL provide the following guarantees:

\begin{description}[style=nextline, leftmargin=3.5em, labelwidth=3em]

  \item[\textbf{G-1}] \textbf{Snapshot integrity.}\enspace
    Any modification to any file within a signed snapshot SHALL be
    detectable by a verifier with access to the public key and the
    snapshot's \texttt{SIGNATURE.json}.  This goal mitigates \textbf{A-1}
    and \textbf{A-2}.

  \item[\textbf{G-2}] \textbf{Deterministic reproducibility.}\enspace
    Any conforming implementation executing the pipeline on identical inputs
    and methodology parameters SHALL produce byte-identical output, including
    identical hash digests and an identical signature preimage.
    This goal mitigates \textbf{A-5}.

  \item[\textbf{G-3}] \textbf{Tamper detectability.}\enspace
    Verification of a snapshot SHALL require only the snapshot directory,
    the public key, and the verification algorithm.  No network access,
    database query, or access to the computation environment SHALL be
    required.  A failed verification SHALL produce an unambiguous rejection.

  \item[\textbf{G-4}] \textbf{Non-repudiation.}\enspace
    A validly signed snapshot constitutes a binding commitment by the
    custodian (\cref{sec:gov:custodian}) to the scores, classifications,
    and rankings contained therein.  The custodian SHALL NOT deny
    authorship of a snapshot bearing a valid signature under a key they
    have published.

\end{description}

% -------------------------------------------------------------------------
\subsection{Out-of-Scope Threats}\label{sec:threat:out-of-scope}

The following threats are explicitly excluded from the security model:

\begin{enumerate}
  \item \textbf{Raw data falsification.}\enspace  The integrity chain
    authenticates the \emph{computation}, not the veracity of the
    underlying source data.  If a statistical agency publishes erroneous
    trade data, the \ISI pipeline will faithfully compute and sign a
    snapshot reflecting that error.  Source-data validation is the
    responsibility of the data providers, not of the \ISI integrity system.

  \item \textbf{Upstream statistical bias.}\enspace  Systematic biases in
    source-data collection (e.g., under-reporting of certain trade flows)
    are not detectable or correctable by the integrity chain.

  \item \textbf{Compromised data providers.}\enspace  If a data provider is
    coerced or incentivised to publish misleading data, the \ISI system has
    no mechanism to detect or correct the resulting distortion.  The
    integrity chain guarantees that the published snapshot correctly reflects
    the computation; it does not guarantee that the inputs to the
    computation are truthful.

  \item \textbf{Side-channel attacks on the computation environment.}\enspace
    Timing attacks, power analysis, and electromagnetic emanation attacks
    against the hardware executing the pipeline are not considered.  The
    pipeline is not a cryptographic primitive and does not process secret
    inputs.

  \item \textbf{Denial of service.}\enspace  Availability of the API
    serving layer is an operational concern governed by infrastructure
    policy, not by this specification.  Rate limiting
    (\cref{sec:api:rate-limiting}) provides a baseline defence but does not
    constitute a formal availability guarantee.
\end{enumerate}

% -------------------------------------------------------------------------
\subsection{Trust Anchors}\label{sec:threat:trust-anchors}

The security model rests on the following trust assumptions.  Violation of
any assumption invalidates the corresponding security goal.

\begin{description}[style=nextline, leftmargin=3.5em, labelwidth=3em]

  \item[\textbf{TA-1}] \textbf{Public key distribution.}\enspace
    The Ed25519 public key is distributed through an authenticated channel
    (e.g., the project repository, a TLS-protected webpage, or a published
    document bearing an institutional identifier).  Verifiers MUST obtain
    the public key from a trusted source.  If the public key is substituted,
    \textbf{G-1} and \textbf{G-4} are void.

  \item[\textbf{TA-2}] \textbf{Append-only registry.}\enspace
    The snapshot's methodology-registry file is append-only at the
    deployment level.  Existing entries SHALL NOT be
    modified or deleted.  New entries MAY be appended for new vintages or
    methodology versions.  This assumption underpins \textbf{G-1}.

  \item[\textbf{TA-3}] \textbf{Offline verification path.}\enspace
    A verifier with access to the snapshot directory and the public key
    SHALL be able to perform full integrity verification without network
    connectivity.  The verification algorithm is specified in
    \\cref{sec:integrity:verify} and is reproducible from this
    document alone.

  \item[\textbf{TA-4}] \textbf{Signing-key confidentiality.}\enspace
    The Ed25519 private key is held exclusively by the custodian
    (\cref{sec:gov:custodian}) and is not shared with data providers,
    API operators, or third parties.  This assumption underpins
    \textbf{G-4}.

\end{description}

% -------------------------------------------------------------------------
\subsection{Mitigation Matrix}\label{sec:threat:mitigation}

\Cref{tab:threat-mitigation} maps each adversary class to the integrity
mechanism that mitigates it and the security goal that is preserved.

\begin{table}[htbp]
  \centering\small
  \caption{Adversary-class mitigation matrix.}\label{tab:threat-mitigation}
  \begin{tabularx}{\textwidth}{@{}l l X l@{}}
    \toprule
    \textbf{Adversary} & \textbf{Mechanism} & \textbf{Detection method} &
      \textbf{Goal} \\
    \midrule
    \textbf{A-1} & SHA-256 + Ed25519 & Hash mismatch on tampered file;
      signature verification failure. & \textbf{G-1} \\
    \textbf{A-2} & Manifest binding (\textbf{C-3}) & Manifest hash covers
      all files; substituted file produces hash mismatch. & \textbf{G-1} \\
    \textbf{A-3} & ETag + offline verification & Client recomputes hashes
      from downloaded snapshot; TLS protects transport. & \textbf{G-3} \\
    \textbf{A-4} & Registry in manifest & Registry file is hashed in
      manifest; mutation invalidates snapshot hash. & \textbf{G-1} \\
    \textbf{A-5} & Determinism invariants (D-1--D-10) & Golden-file and
      cross-platform tests detect drift. & \textbf{G-2} \\
    \textbf{A-6} & Key rotation protocol & Revocation + re-signing with
      new key; see \cref{sec:threat:key-rotation}. & \textbf{G-4} \\
    \bottomrule
  \end{tabularx}
\end{table}

% -------------------------------------------------------------------------
\subsection{Key Rotation and Revocation Policy}\label{sec:threat:key-rotation}

The following normative rules govern the Ed25519 signing key lifecycle:

\begin{enumerate}
  \item The custodian SHALL generate a new Ed25519 key pair when the
    methodology undergoes a major-version increment.

  \item The custodian MAY generate a new key pair at any time for
    operational reasons (e.g., personnel change, infrastructure migration).

  \item Upon key rotation, the custodian SHALL publish the new public key
    through the same authenticated channels used for the previous key
    (\textbf{TA-1}).

  \item The previous public key SHALL remain published and labelled as
    ``superseded'' with an effective date.  Removal of a superseded public
    key is prohibited: verifiers of historical snapshots require access to
    the key that was active at the time of signing.

  \item Snapshots signed with a superseded key SHALL NOT be re-signed with
    the new key.  Each snapshot's signature is bound to the key that was
    active at the time of publication.  Re-signing would violate the
    immutability contract (\cref{sec:snapshot:immutability}).

  \item In the event of suspected or confirmed key compromise
    (\textbf{A-6}), the custodian SHALL:
    \begin{enumerate}[label=(\alph*)]
      \item Immediately revoke the compromised key by publishing a signed
        revocation notice using the compromised key (if still in possession)
        or via an authenticated institutional channel.
      \item Generate and publish a new key pair.
      \item Audit all snapshots signed with the compromised key and publish
        an integrity report identifying any snapshots whose authenticity
        cannot be confirmed.
      \item Re-sign affected snapshots with the new key only if the
        custodian can independently verify that the snapshot contents are
        authentic.  Re-signed snapshots are published in an addendum
        directory alongside the original (\cref{sec:snapshot:immutability}).
    \end{enumerate}

  \item Key pairs SHALL be generated using a cryptographically secure
    random number generator conforming to the requirements of
    \textcite{rfc8032}.  The private key SHALL be stored in an
    access-controlled environment with audit logging.
\end{enumerate}
